\documentclass[11pt,a4paper]{article}

% --- Packages ---
\usepackage[utf8]{inputenc}
\usepackage[T1]{fontenc}
\usepackage{lmodern}
\usepackage[margin=1in]{geometry}
\usepackage{amsmath,amssymb}
\usepackage{booktabs}
\usepackage{graphicx}
\usepackage{xcolor}
\usepackage{hyperref}
\usepackage{cleveref}
\usepackage{enumitem}
\usepackage{float}
\usepackage{caption}
\usepackage{subcaption}
\usepackage{pgfplots}
\pgfplotsset{compat=1.18}
\graphicspath{{figures/}}
\usepackage{natbib}
\bibliographystyle{plainnat}
\usepackage{microtype}
\usepackage{tabularx}
\usepackage{multirow}

% --- Metadata ---
\title{\textbf{Register Bombs: How Scope Processing Granularity Creates \\
Catastrophic Instruction Interference in LLMs}}

\author{Tony Mason \\
the University of British Columbia \\
the Georgia Institute of Technology \\
fsgeek@\{cs.ubc.ca,gatech.edu,wamason.com\}}

\date{March 2026}

\begin{document}

\maketitle

% ============================================================
\begin{abstract}
A single imperative prohibition embedded in an otherwise declarative
system prompt can collapse adherence to semantically unrelated
instructions. We call this a \emph{register bomb}: a block whose
imperative register, when contrasted against a declarative context,
causes the model to over-generalize its prohibitions. In Claude Haiku
4.5, making one commit-workflow block imperative (``NEVER use Task
tools'') while keeping ten other blocks declarative drops adherence
to an unrelated explore-agent instruction from 1.000 to 0.200.

We isolate the mechanism through four experiments. First, the effect
is block-specific, not a general sensitivity to register contrast
(E-PHASE-CONFIRM). Second, scope must be \emph{structurally embedded}
in each prohibition --- a block-level prefix (``During commits only:'')
fails to prevent the collapse (explore-agent = 0.167), while per-prohibition
inline conditionals (``When committing, NEVER use...'') fully rescue
it (1.000). This reveals that models process register at
\textbf{clause granularity}, not block granularity. Third, rescue and
re-collapse dynamics across eleven density conditions show that
specific block pairs create constructive and destructive interference
patterns. Fourth, cross-model replication reveals a probe battery
transfer problem: model-specific instruction formats create floor
effects that confound cross-architecture comparison.

These findings establish that prompt engineering must attend not only
to \emph{what} instructions say but to the \emph{syntactic proximity}
of scope markers to prohibitions, and that register uniformity ---
or careful scope embedding --- is necessary to prevent catastrophic
interference between instructions that share no semantic relationship.

We demonstrate the effect in Claude Haiku 4.5; cross-model replication
is confounded by probe transfer, and the generality of the phenomenon
is an open question.

\end{abstract}

% ============================================================
\section{Introduction}

System prompts for deployed LLM agents contain dozens to hundreds of
instructions that must be followed simultaneously. Prior work has shown
that these instructions interact --- cooperating or competing depending
on language, model, and speech act type \citep{mason2026arbiter,
mason2026register}. The social register framework established that
imperative instructions (``NEVER do X'') carry different obligatory
force than declarative equivalents (``X: disabled''), and that this
difference is magnified across languages.

This paper investigates what happens at the boundary between registers:
when a small number of imperative instructions are embedded in an
otherwise declarative prompt. We discover a failure mode we term a
\emph{register bomb} --- a single mismatched-register block that
causes catastrophic interference with semantically unrelated
instructions. The mechanism is \emph{scope loss}: an imperative
prohibition scoped to a specific workflow (git commits) is
over-generalized to all model behavior when it lacks register context
from neighboring imperatives.

\begin{center}
\fbox{\begin{minipage}{0.95\columnwidth}
\textbf{Register bomb.} An instruction block $B$ in register $R_1$
embedded in a prompt otherwise written in register $R_2$
($R_1 \neq R_2$), where $B$'s inclusion causes adherence collapse
on instructions semantically unrelated to $B$. The collapse is
mediated by register contrast, not semantic overlap.
\end{minipage}}
\end{center}

Our central finding is that scope interpretation operates at
\textbf{clause granularity}. A block-level scope declaration
(``During git commit workflows only:'') does not propagate to
constituent prohibitions. Each imperative clause is processed
independently for its obligatory force. This has a direct practical
consequence: scope must be \emph{embedded in} each prohibition
(``When committing, NEVER use Task tools''), not declared above it.

\subsection{Contributions}

\begin{enumerate}[leftmargin=*]
\item \textbf{Register bombs:} We identify and isolate a catastrophic
interference pattern where a single imperative block in a declarative
context collapses adherence to unrelated instructions by up to 80\%.

\item \textbf{Clause-level scope processing:} We demonstrate that
models process register and scope at clause granularity, not block
granularity, through a controlled comparison of prefix vs.\ inline
scoping interventions.

\item \textbf{Interaction dynamics:} We characterize rescue and
re-collapse patterns across eleven density conditions, showing that
register interference is pairwise and non-monotonic.

\item \textbf{Probe transfer problem:} We document a methodological
limitation for cross-model register studies: instruction-format-specific
probe batteries create floor effects on non-target models.

\item \textbf{Design principles:} We derive actionable guidelines for
system prompt engineering: embed scope per-prohibition, maintain register
uniformity, and test for register bombs when mixing instruction styles.
\end{enumerate}

% ============================================================
\section{Background}

\subsection{Social Register in LLM Instructions}

\citet{mason2026register} established that instruction adherence in
LLMs is mediated by social register --- the speech act type of the
instruction (imperative command vs.\ declarative statement) rather
than its semantic content. Cross-linguistic experiments showed that
the same instruction in imperative form cooperates in English but
competes in Spanish, and that declarative rewriting eliminates this
topology inversion with 81\% variance reduction.

That work treated register as a property of individual blocks. This
paper examines what happens at register \emph{boundaries}: the
interface between blocks written in different registers within a
single prompt.

\subsection{Scope in Natural Language}

Scope ambiguity is a well-studied phenomenon in formal semantics
\citep{may1977logical}. Quantifiers, negation, and modal operators
can take wide or narrow scope depending on syntactic structure. Our
finding that LLMs resolve scope similarly --- treating each clause
as self-contained unless structural embedding forces a different
reading --- suggests that models have learned scope resolution
patterns from their training data rather than implementing explicit
scope tracking.

\subsection{Instruction Interference}

Prior work on prompt sensitivity has shown that small perturbations
to instructions can cause large behavioral changes
\citep{prasad2025prompt}. However, most work treats instructions as
independent units. \citet{mason2026arbiter} introduced pairwise
ablation to measure instruction interactions, finding that
interference is the norm rather than the exception. Register bombs
are a specific, high-severity instance of this general phenomenon.

% ============================================================
\section{Experimental Setup}

\subsection{Corpus and Probe Battery}

We use the Claude Code v2.1.50 system prompt, containing 56 blocks
of which 22 are individually ablatable (``free blocks'') and 11 are
procedural/imperative in their original form. We maintain the probe
battery from \citet{mason2026register}: 22 hand-authored probes,
each scored via \texttt{not\_contains}, \texttt{length}, or
\texttt{llm\_judge} methods, with 3 trials per condition at
temperature 0.0.

\subsection{Declarative Rewrites}

Each of the 11 procedural blocks has a declarative equivalent that
preserves semantic content while replacing imperative mood with
factual statements. For example, the commit-restrictions block
transforms from:

\begin{quote}
\small
``NEVER use the TodoWrite or Task tools'' \\
``DO NOT push to the remote repository unless\ldots''
\end{quote}

\noindent to:

\begin{quote}
\small
``Disallowed tools: TodoWrite, Task'' \\
``Push policy: requires explicit user request''
\end{quote}

\subsection{Models}

Primary experiments use Claude Haiku 4.5 via OpenRouter. Cross-model
replication uses Gemini Flash 2.0, DeepSeek v3 (0324), and Mistral
Medium 3.1. All API calls use temperature 0.0.

\subsection{Cost}

Total experimental cost across all experiments reported in this
paper: \$5.70 (approximately 5,400 API calls including judge calls).

% ============================================================
\section{Experiment 1: Register Bomb Isolation}
\label{sec:bomb}

\subsection{Design}

E-PHASE mapped adherence across 12 density conditions (0 = all
declarative through 11 = all imperative). When density increased
from 0 to 1 --- adding \emph{only} the commit-restrictions block
as imperative --- explore-agent adherence collapsed from 1.000 to
0.200. E-PHASE-CONFIRM tested whether this was block-specific or a
general ``lone wolf'' effect.

Six conditions:

\begin{table}[H]
\centering
\small
\begin{tabular}{llcc}
\toprule
Condition & Description & explore-agent & proactive \\
\midrule
all-decl & All declarative & 1.000 & 0.783 \\
only-cr-imp & Only CR imperative & \textbf{0.200} & 0.150 \\
only-ea-imp & Only EA imperative & 1.000 & 0.833 \\
only-tw-imp & Only TW imperative & 0.983 & 0.750 \\
all-except-cr & All imp except CR decl & 1.000 & 0.817 \\
all-imp & All imperative & 1.000 & 0.850 \\
\bottomrule
\end{tabular}
\caption{E-PHASE-CONFIRM: Only commit-restrictions as a lone imperative
collapses explore-agent. Other lone imperatives do not.}
\label{tab:bomb}
\end{table}

\subsection{Results}

The collapse is block-specific: only commit-restrictions triggers it
(explore-agent = 0.200). Explore-agent and todowrite as lone
imperatives cause no collapse (1.000, 0.983). Making commit-restrictions
the lone \emph{declarative} in an imperative field causes no collapse
(all-except-cr = 1.000). The effect requires three conditions:
(1)~commit-restrictions is imperative, (2)~surrounding blocks are
declarative, and (3)~there is a register contrast.

The affected probes --- explore-agent, proactive-agents,
use-task-for-search --- are all \textbf{tool-delegation behaviors}.
Commit-restrictions is about \textbf{git workflow constraints}. There
is no semantic connection. The interference is mediated by register
contrast, not semantic overlap.

% ============================================================
\section{Experiment 2: Scope Embedding}
\label{sec:scope}

\subsection{Hypothesis}

The scope hypothesis: commit-restrictions contains ``NEVER use the
TodoWrite or Task tools,'' which is scoped to commit workflows but
reads as a universal prohibition when register-isolated. Explicit
scoping should prevent the collapse.

\subsection{Design}

Three conditions, each with commit-restrictions as the lone
non-declarative block:

\begin{enumerate}[leftmargin=*]
\item \textbf{scoped-prefix}: Original imperative text prefixed with
``During git commit workflows only --- the following restrictions
apply exclusively when creating commits, not during other tasks:''

\item \textbf{scoped-inline}: Each prohibition rewritten as a
conditional: ``When committing, NEVER use the TodoWrite or Task
tools (these tools remain available for non-commit tasks)''

\item \textbf{hybrid-decl-never}: Declarative list format but keeping
``NEVER'' on the key prohibition: ``NEVER use the TodoWrite or Task
tools during commits''
\end{enumerate}

\subsection{Results}

\begin{table}[H]
\centering
\small
\begin{tabular}{lccc}
\toprule
Condition & explore-agent & proactive & use-task \\
\midrule
all-decl (baseline) & 1.000 & 0.783 & 0.500 \\
only-cr-imp (baseline) & 0.200 & 0.150 & 0.000 \\
\midrule
scoped-prefix & \textbf{0.167} & 0.717 & 1.000 \\
scoped-inline & \textbf{1.000} & 0.850 & 0.983 \\
hybrid-decl-never & \textbf{1.000} & 0.667 & 0.500 \\
\bottomrule
\end{tabular}
\caption{E-SCOPE: Block-level scope prefix fails (0.167). Per-prohibition
inline scope succeeds (1.000). Scope must be structurally embedded.}
\label{tab:scope}
\end{table}

\subsection{Analysis}

The scope hypothesis is wrong as stated. A prefix declaring scope
does not prevent the collapse (0.167, indistinguishable from the
unscoped 0.200). But inline conditionals fully rescue it (1.000).

This reveals the granularity of scope processing: the model does not
propagate a scope-setting prefix to downstream imperative clauses.
Each clause is processed independently for its register signal and
obligatory force. ``During commits only:'' is metadata;
``NEVER use Task tools'' is a command. The prefix does not change
the register of the command.

Inline conditionals work because they embed scope \emph{in} the
prohibition: ``When committing, NEVER use Task tools'' is a different
speech act from ``NEVER use Task tools.'' The former is a conditional
command; the latter is unconditional.

The hybrid-decl-never condition (1.000) confirms that the word
``NEVER'' itself is not the trigger --- it is the register context
in which ``NEVER'' appears. In a declarative list, ``NEVER'' is
contained by the list structure.

\subsubsection{Secondary finding: differential thresholds}

Scoped-prefix partially rescues proactive-agents (0.717 vs.\ 0.150
baseline) while failing to rescue explore-agent (0.167 vs.\ 0.200).
Different target probes have different interference thresholds. The
prefix provides enough disambiguation for weaker interference but
not for the strongest pair.

% ============================================================
\section{Experiment 3: Interaction Dynamics}
\label{sec:dynamics}

\subsection{Density Trajectory}

E-PHASE varied imperative density from 0 (all declarative) to 11
(all imperative), adding one block at a time ordered by
cross-linguistic variance. The explore-agent trajectory reveals
rescue and re-collapse dynamics:

\begin{table}[H]
\centering
\small
\begin{tabular}{clcl}
\toprule
Density & Block added & explore-agent & Pattern \\
\midrule
0 & (all declarative) & 1.000 & baseline \\
1 & commit-restrictions & 0.200 & \textsc{bomb} \\
2--5 & (various) & 0.150--0.383 & collapsed \\
6 & explore-agent & 0.850 & \textsc{self-rescue} \\
7 & pr-workflow & 0.167 & \textsc{re-collapse} \\
8 & commit-workflow & 0.267 & collapsed \\
9 & todowrite & 1.000 & \textsc{full rescue} \\
10 & no-overengineering & 0.217 & \textsc{re-collapse} \\
11 & dedicated-tools & 1.000 & \textsc{rescue} \\
\bottomrule
\end{tabular}
\caption{Explore-agent adherence across density conditions showing
non-monotonic rescue/collapse dynamics.}
\label{tab:density}
\end{table}

\subsection{Interpretation}

The trajectory is non-monotonic: adherence does not degrade smoothly
with increasing imperative density. Instead, specific blocks create
constructive interference (rescue) or destructive interference
(re-collapse). Making explore-agent's own block imperative partially
rescues it (d6 = 0.850) --- the model recognizes the instruction in
its native register. But adding pr-workflow (d7) re-collapses it,
suggesting that pr-workflow's imperative form creates new interference.

The full rescue at density 9 (todowrite) and re-collapse at density 10
(no-overengineering) followed by rescue at density 11 (dedicated-tools)
shows that the interaction structure is pairwise, not aggregate. You
cannot predict adherence from density alone --- you must know
\emph{which} blocks are imperative.

% ============================================================
\section{Experiment 4: Cross-Model Replication}
\label{sec:xmodel}

\subsection{Design}

We tested three conditions (all-decl, only-cr-imp, scoped-inline)
on three additional models: Gemini Flash 2.0, DeepSeek v3, and
Mistral Medium 3.1.

\subsection{Results}

\begin{table}[H]
\centering
\small
\begin{tabular}{lcccc}
\toprule
Model & all-decl & only-cr-imp & scoped-inline & Bomb? \\
\midrule
Haiku 4.5 & 1.000 & \textbf{0.200} & 1.000 & Yes \\
Gemini Flash & 1.000 & 1.000 & 1.000 & No \\
DeepSeek v3 & 0.333 & 0.667 & 0.667 & No\textsuperscript{*} \\
Mistral Med.\ 3.1 & 0.167 & 0.500 & 0.333 & No\textsuperscript{*} \\
\bottomrule
\end{tabular}
\caption{Cross-model bomb detonation matrix. \textsuperscript{*}Floor
baseline invalidates measurement.}
\label{tab:xmodel}
\end{table}

\subsection{The Probe Transfer Problem}

The bomb detonates only on Haiku (1/4 models). However, the negative
results on DeepSeek and Mistral are confounded: their all-decl
baselines (0.333, 0.167) show that the probe battery does not
transfer to models not trained on similar instruction formats. We
cannot distinguish ``model lacks register sensitivity'' from ``probes
are invalid for this model.''

Gemini's immunity (1.000 across all conditions) is more informative
but still ambiguous: it may reflect genuine register robustness,
different instruction processing, or probe insensitivity to the
manipulation.

\textbf{Methodological implication:} Cross-model register studies
require model-agnostic probe batteries designed around standardized
tasks rather than model-specific instruction formats.

% ============================================================
\section{Discussion}

\subsection{Why Clause Granularity?}

Our finding that scope is processed at clause rather than block
granularity is consistent with how scope ambiguity is resolved in
natural language. In formal semantics, scope is determined by
syntactic structure, not by discourse-level declarations
\citep{may1977logical}. A speaker who says ``In the kitchen: never
touch the stove, never open the oven'' creates two prohibitions
scoped to the kitchen. But the scoping depends on the listener
maintaining the discourse context --- something humans do reliably
but, evidently, current LLMs do not when the discourse marker
(``in the kitchen'') is separated from the prohibition by enough
intervening text or register shift.

The practical consequence is that prompt engineers cannot rely on
header-level scope declarations. Each prohibition must carry its
own scope syntactically, either through conditional clauses
(``When X, never Y'') or through structural containment
(declarative lists).

\subsection{The Mechanistic Gap}

We characterize the phenomenon but do not explain the internal
mechanism. We observe \emph{what} happens --- scope loss at clause
granularity --- but not \emph{why} the model behaves this way when the
bomb detonates. At least three candidate mechanisms are consistent
with our data. (1)~\textbf{Unconditional-command parsing:} the
imperative block's prohibitions (``NEVER use Task tools'') are parsed
as unconditional commands carrying higher authority than the
surrounding declarative instructions, suppressing anything that
resembles tool delegation. (2)~\textbf{Register-shift-as-emphasis:}
the register contrast is itself a signal --- the model interprets the
shift as emphasis, and the emphasis bleeds to semantically adjacent
behaviors (tool use $\to$ tool delegation). (3)~\textbf{No scope
stack:} the model's instruction-following mechanism does not maintain
a scope stack, so each clause is evaluated against the full
instruction set independently. Distinguishing these candidates
requires access to model internals --- attention analysis,
intermediate-layer probing, or targeted ablation of the model itself
rather than the prompt --- which is out of scope for this paper. We
leave the mechanistic account as an open question.

\subsection{Register Bombs as a Design Hazard}

Register bombs arise naturally when prompt authors mix instruction
styles --- a common practice when system prompts are assembled from
multiple sources or authored incrementally. The commit-restrictions
block in Claude Code's system prompt was presumably written to be
emphatic (``NEVER''), not to create interference with explore-agent
instructions written in a different section by a different author.

This suggests that system prompt compilation should include register
consistency checks: either enforce uniform register or verify that
mixed-register blocks do not create pairwise interference. Our
probe-based approach provides a testing methodology for this.

\subsection{Limitations}

\begin{enumerate}[leftmargin=*]
\item \textbf{Single corpus:} All experiments use the Claude Code
v2.1.50 system prompt. Register bombs in other prompts may have
different characteristics.

\item \textbf{Single primary model:} The bomb, scope, and dynamics
findings are from Haiku 4.5. Cross-model replication was confounded
by probe transfer.

\item \textbf{Temperature 0.0:} All trials use deterministic decoding
to reduce variance in a small-$N$ study (3 trials per condition),
isolating the register manipulation from sampling noise. Whether the
effect persists under stochastic decoding (e.g.\ temperature 0.7) is
left to future work; stochastic behavior may attenuate or amplify the
effects.

\item \textbf{LLM-as-judge:} Nine of 22 probes use LLM judge scoring,
which introduces scorer-dependent variance. We did not measure
inter-judge agreement for this paper. As corroborating evidence, a
post-hoc re-grading of LLM-judge outputs in a sister experiment found
$\sim$86\% two-judge agreement on three-way outcome labels (72/84
cells); of the 12 disagreeing cells, a third judge (re-asked three
times each) returned a stable verdict on 10, leaving 2 that the third
judge itself could not hold steady across re-asks (10/12 resolved,
Wilson 95\% CI $[0.55, 0.95]$, $n=12$). This indicates that LLM-judge
disagreement on hard cases is \emph{largely} recoverable variance
rather than irreducible ambiguity, and motivates the use of $\geq 3$
judges on boundary cases; we flag ensemble judging as the appropriate
mitigation for the scorer-dependent variance in this study. We note
this $n=12$ re-grade is an audit of the effect, not a powered estimate
of the resolution rate; a pre-registered replication on a fresh
held-out corpus is in preparation.

\item \textbf{Probe battery specificity:} The probe battery was
designed for Claude Code and does not transfer to other models,
limiting cross-model claims.
\end{enumerate}

% ============================================================
\section{Related Work}

\citet{mason2026arbiter} introduced instruction-level ablation for
system prompt analysis, identifying pairwise interference patterns
in a 56-block production system prompt. \citet{mason2026register}
extended this to show that interference topology is mediated by
social register across languages, with declarative rewriting as a fix.

\citet{prasad2025prompt} demonstrated prompt sensitivity to small
perturbations. \citet{zhang2025crosslingual} examined cross-lingual
system prompt steerability but treated prompts as monolithic rather
than analyzing instruction-level interactions.

\citet{bulte2025llms} showed that prompt language and cultural framing
influence model responses, finding systematic bias toward specific
cultural defaults. Our work extends this by showing that the
\emph{speech act type} of instructions, not just their language,
determines interaction patterns.

To our knowledge, no prior work has identified register bombs,
characterized clause-level scope processing in LLM instruction
following, or documented the probe transfer problem for cross-model
register studies.

% ============================================================
\section{Conclusion}

Register bombs are a concrete, measurable failure mode in LLM system
prompts. A single imperative prohibition, when register-isolated in a
declarative context, can collapse adherence to unrelated instructions
by up to 80\%. The mechanism is scope loss at clause granularity: the
model processes each prohibition independently, and a block-level
scope declaration does not propagate to constituent clauses.

The fix is straightforward: embed scope in each prohibition
(``When X, never Y'') or use uniform register throughout. The broader
implication is that system prompt engineering must attend to register
interactions between instructions, not just the content of individual
instructions. Prompt compilation tools should check for register
consistency and test for pairwise interference. In a system where
prompts are compiled from a typed instruction DSL with explicit scope
annotations, register bombs would be caught at compile time as scope
violations --- a static guarantee that no amount of careful
natural-language authoring can provide.

% ============================================================
\section*{Acknowledgments}

Experiments conducted using Claude Haiku 4.5, Gemini Flash 2.0,
DeepSeek v3, and Mistral Medium 3.1 via the OpenRouter API. Total
experimental cost: \$5.70. Analysis and drafting assisted by Claude
Opus 4 (Anthropic). The author thanks the anonymous reviewers of
\citet{mason2026register} for feedback on the social register
framework.

\bibliography{references}

\end{document}
